\chapter{Acne}
\index{Acne}

\section*{Definition and Overview}
Acne is a very common skin condition that occurs when hair follicles become blocked with oil (sebum) and dead skin cells, leading to blackheads, whiteheads, pimples, and, in more severe cases, deeper nodules and cysts. It most often affects the face but also commonly appears on the back, chest, and shoulders, which are rich in oil-producing glands.

Acne is not a sign of poor hygiene and it is not infectious. Although it is usually seen in teenagers, it can start or persist in adulthood, and it can cause significant distress and permanent scarring. A range of effective treatments exists, from over-the-counter creams to specialist prescription medicines, and most cases can be controlled well.

\section*{Epidemiology}
Acne is the most common skin condition worldwide, affecting most people at some point in their lives. Roughly 8 in 10 people aged 11 to 30 experience acne at least occasionally, and it is most common during the teenage years when hormone levels change.

Males are more often severely affected during adolescence, whereas acne that continues or first appears in adulthood is more common in women. Acne affects people of all ethnicities, and many women experience flare-ups around the menstrual cycle. It is also more common in certain conditions associated with elevated androgens, such as polycystic ovary syndrome.

\section*{Aetiology and Pathophysiology}
Acne develops when four factors act together in the pilosebaceous unit, the hair follicle and its associated oil gland.

\begin{itemize}
    \item \textbf{Excess sebum:} hormonal changes, particularly increased androgens, stimulate the oil glands to produce more oil
    \item \textbf{Blocked follicles:} the opening of the hair follicle becomes blocked by a plug of oil and dead skin cells
    \item \textbf{Bacteria:} the skin bacterium \textit{Cutibacterium acnes} multiplies inside the blocked follicle
    \item \textbf{Inflammation:} the contents of the follicle spill into the surrounding skin, producing red, swollen spots
\end{itemize}

Genetics play a large role in determining who gets acne and how severely. Hormonal changes at puberty, with the menstrual cycle, and in pregnancy can trigger flare-ups. Contrary to popular belief, chocolate and greasy foods are not proven causes, and dirt does not cause acne; however, oily cosmetics, heavy sweating, friction from clothing or helmets, and some medicines (including corticosteroids and certain epilepsy treatments) can make it worse.

\section*{Clinical Features}
\begin{itemize}
    \item \textbf{Blackheads (open comedones):} small dark dots caused by a blocked follicle opening
    \item \textbf{Whiteheads (closed comedones):} small flesh-coloured or white bumps
    \item \textbf{Papules:} small, red, tender bumps
    \item \textbf{Pustules:} red bumps with a white or yellow centre of pus
    \item \textbf{Nodules and cysts:} larger, deeper, painful lumps beneath the skin
    \item \textbf{Distribution:} a greasy complexion, with spots most often on the face, back, chest, and shoulders
    \item \textbf{Residual change:} dark marks (post-inflammatory hyperpigmentation) or permanent scars may remain after spots heal
\end{itemize}

\section*{Differential Diagnosis}
A number of skin conditions can be mistaken for acne and should be distinguished.

\begin{itemize}
    \item Rosacea (adult facial redness and pimples, usually without blackheads)
    \item Folliculitis (irritation or infection of the hair follicles)
    \item Perioral dermatitis (small red bumps around the mouth and nose)
    \item Hidradenitis suppurativa (boil-like lumps in the armpits and groin)
    \item Keratosis pilaris (rough bumps on the upper arms)
    \item Milia (tiny white cysts, common in infants)
    \item Contact dermatitis or bacterial skin infections
\end{itemize}

\section*{When to Seek Medical Care}
See a doctor if acne is severe, painful, or leaving scars, if over-the-counter treatments have not helped after several months, or if it is causing significant distress. For women, acne accompanied by irregular periods, excessive facial hair, or difficulty losing weight may point to a hormonal condition that needs investigation.

\textbf{Contact your local emergency services for:}
\begin{itemize}
    \item A high fever with a painful, rapidly spreading skin infection
    \item A large, swollen, and increasingly painful area of skin that is spreading
    \item Sudden confusion, severe drowsiness, or difficulty breathing
    \item Any sudden, severe, or concerning symptom
\end{itemize}

\section*{Diagnosis and Investigations}
\begin{itemize}
    \item \textbf{History and examination:} the diagnosis is usually made by looking at the skin and asking about medicines, cosmetics, and menstrual history
    \item \textbf{Severity assessment:} classifying the types and number of spots helps guide treatment choice
    \item \textbf{Blood tests:} may be arranged for women with signs of a hormonal problem, such as irregular periods or excessive hair growth
    \item \textbf{Microbiology:} a swab is occasionally taken if a bacterial infection is suspected
    \item \textbf{Specialist referral:} a dermatologist may be involved for severe, scarring, or treatment-resistant acne
\end{itemize}

\section*{Management}
\begin{itemize}
    \item \textbf{Gentle skin care:} washing twice daily with a mild cleanser and using non-comedogenic (non-pore-blocking) products
    \item \textbf{Over-the-counter products:} treatments containing benzoyl peroxide, salicylic acid, or sulphur
    \item \textbf{Topical prescription medicines:} retinoids (vitamin A derivatives) and antibiotic creams or gels
    \item \textbf{Oral medicines:} antibiotic tablets for a limited course, and hormonal treatments such as the combined oral contraceptive pill for some women
    \item \textbf{Isotretinoin:} a powerful vitamin A medicine for severe or scarring acne, supervised by a specialist because of its side effects
    \item \textbf{Procedures:} extraction of blackheads, chemical peels, and light or laser therapies in selected cases
    \item \textbf{Follow-up:} most treatments take months to work, so treatment is usually reviewed after 2 to 3 months
\end{itemize}

\section*{Complications}
\begin{itemize}
    \item Permanent scarring and dark marks, which can persist after the acne has settled
    \item Psychological effects, including low self-esteem, anxiety, and depression
    \item Social withdrawal and reduced participation in school, work, and sport
    \item Side effects of treatment, such as skin irritation from topical products or effects of oral medicines
    \item The specific risks of isotretinoin, which requires careful supervision and cannot be used in pregnancy
\end{itemize}

\section*{Prognosis and Outcomes}
Acne usually improves with age, and most people's skin clears during their late teens or twenties. With treatment, most cases can be controlled well, although 2 to 3 months of regular treatment is commonly needed before a clear improvement is seen.

Some people have acne that persists or recurs into adulthood and needs longer-term management. Severe acne can cause permanent scarring, but prompt treatment reduces this risk, and a course of isotretinoin can lead to long-lasting clearance in many people. Consistent self-care and adherence to treatment give the best chance of good skin and few scars.

\section*{Prevention and Self-Care}
\begin{itemize}
    \item Wash the skin twice daily with a gentle cleanser and warm water; avoid harsh scrubbing
    \item Do not pick, squeeze, or pop spots, as this increases scarring and infection
    \item Use non-comedogenic, oil-free moisturisers, sunscreens, and cosmetics
    \item Wash the face after heavy sweating and change pillowcases regularly
    \item Manage stress and get adequate sleep, which can influence flare-ups
    \item Seek treatment early for moderate or severe acne rather than waiting for it to clear
\end{itemize}

\section*{Sources and Further Reading}
The following organisations provide reliable information on acne and its treatment.

\begin{itemize}
    \item NHS. \textit{Information on acne, including causes and treatments.} https://www.nhs.uk/conditions/
    \item Mayo Clinic. \textit{Guide to acne, diagnosis, and management options.} https://www.mayoclinic.org/
    \item MedlinePlus. \textit{Health information on acne and skin conditions.} https://medlineplus.gov/
    \item MSD Manual. \textit{Professional and patient education on acne.} https://www.msdmanuals.com/
    \item World Health Organization. \textit{Resources on skin health and adolescent health.} https://www.who.int/
\end{itemize}
